\documentclass[12pt, a4paper,twoside]{report}

%% Every LaTeX document begins with a preamble, which loads packages and
%% defines various settings to make the document look right. Mostly,
%% you can ignore everything in this template before \begin{document} on
%% line 74

\usepackage{mathtools,amsthm} % Enable useful mathematical symbols/environments
\usepackage{graphicx} % Enable graphics
\usepackage{fancyhdr,titlesec,microtype} % enable various formatting commands
\usepackage[margin=2cm]{geometry} % Set margin size
\usepackage{palatino} % Set the font
\usepackage[latin1]{inputenc} % Allow you to input accents, umlauts and other characters
\usepackage[T1]{fontenc} % Lets LaTeX print a wider array of characters

\usepackage{xcolor} % Enable coloured elements
\definecolor{mypurple}{HTML}{622567} %%% Purple
\definecolor{myred}{HTML}{D55C19} %%%EssexOrange
\definecolor{myblue}{HTML}{007A87} %%%Seagrass

% For technical reasons, hyperref should be loaded after all other packages
\usepackage[colorlinks,linkcolor=myblue,citecolor=mypurple]{hyperref}

\renewcommand{\baselinestretch}{1.5} % 1.5 line spacing

% Define \begin{theorem}, \end{theorem}, etc.
\theoremstyle{plain} % The following will be italicised
\newtheorem{theorem}{Theorem}[chapter]

\theoremstyle{definition} % The following environments will not use italics
\newtheorem{definition}[theorem]{Definition}
\newtheorem{example}[theorem]{Example}

\numberwithin{equation}{chapter}

% Fancy headings
\pagestyle{fancy}
\setlength{\headheight}{15pt}
\fancyheadoffset[LE,RO]{0pt}
\renewcommand{\chaptermark}[1]{\markboth{#1}{}}
\renewcommand{\sectionmark}[1]{\markright{\thesection\ #1}}
\fancyhf{}
\fancyhead[LE]{\makebox[0pt][l]{\thepage}\hfill\leftmark}
\fancyhead[RO]{\rightmark\hfill\makebox[0pt][r]{\thepage}}
\fancypagestyle{plain}{%
    \fancyhead{} % get rid of headers
    \renewcommand{\headrulewidth}{0pt} % and the line
}

% Fancy chapter numbers
\titleformat{\chapter}[display]
    {\normalfont\bfseries\color{myred}}
    {\filleft\hspace*{-60pt}
        \rotatebox[origin=c]{90}{
            \normalfont\color{black}\large
            \textls[180]{\textsc{\chaptertitlename}}
        }
        \hspace{10pt}
        {\setlength\fboxsep{0pt}
            \colorbox{myred}{\parbox[c][2cm][c]{1.5cm}{
                    \centering\color{white}\fontsize{40}{45}\selectfont\thechapter}
            }
        }
    }
    {10pt}
    {\titlerule[2.5pt]\vskip3pt\titlerule\vskip4pt\LARGE\sffamily}

\begin{document} % Start your document

%%%%%%%%%%%% BEGIN TITLE PAGE %%%%%%%%%%%%

\thispagestyle{empty} % For the title page, no header / footer

\noindent
\begin{minipage}{0.1\textwidth}
    \includegraphics[height=4.5em]{essex.png}
\end{minipage}
\hfill
\begin{minipage}{0.5\textwidth}
    \begin{center}
        \renewcommand\familydefault{\sfdefault}
        \fontfamily{phv}\selectfont
        {\large School of Computer Science and Electrical Engineering}
    \end{center}
\end{minipage}

\begin{center}
    \noindent\textcolor{myred}{\rule{\linewidth}{4.8pt}}

    \vspace{2em}
    \noindent {\LARGE \sc Capstone Project Dissertation}

    \vspace{3em}
    \noindent {\Huge{\color{myblue} GoLearn - Learning Resource Management
            System}}

    \vspace{3em}
    \noindent {\Large \bf Bradley Beasley}

    \vfill
    \noindent {\Large {Supervisor:} {\color{mypurple} \bf Dr Alexandoros
            Voudoris}}

    \vspace{0.5em}
    \noindent\textcolor{myred}{\rule{\linewidth}{4.8pt}}

    \vspace{2em}
    {\Large \today }

    {\Large Colchester}
\end{center}

\clearpage

\abstract{
    There are plenty of services and programs out there for allowing Teachers
    to upload resources and for their learners to access them. Most of these
    are not purpose designed for places of learning, such as school, colleges
    and universities. Those that exist for this specific use-case tend to be
    over-engineered and can be overly complicated (Moodle is a good example
    of this). The purpose of GoLearn is to provide a simple LRM experience
    for teachers and students to interact with the same resources without a
    large amount of overhead in setting up the system: it should work out of
    the box.\\This report contains a look into methods of implementation for
    web applications in general, plus the choices in technology and design
    that were made for GoLearn.
}

%%%%%%%%%%%% END TITLE PAGE %%%%%%%%%%%%

\tableofcontents

%%%%%%%%%%%% Introduction %%%%%%%%%%%%

\chapter{Introduction}\label{ch:introduction}

%%%%%%%%%%%% Background %%%%%%%%%%%%

\chapter{Background}\label{ch:background}

\section{Similar Services}\label{sec:similarservices}

There are various LRMS's available in the market, with some of the more popular
one being:

\begin{description}
    \item[Udemy] {
        is a free to access learning site that puts a paywall over all of
        the courses. Creators can sign up to the site and create a course very
        quickly. Learners can sign up to courses (paying whatever charge is
        applied) and have unlimited access to them. They can access these from
        their own `Courses' dashboard.
    }
    \item[Skillshare] {
        offers a different business model to generate income from
        courses. Where Udemy uses a Course as a Product model, SkillShare
        provides a Subscription service to access all content that they provide.
        Teachers are then paid out based on their course metrics
        \cite{skillshare-teacher-earnings}.\\
        The student experience of SkillShare is very similar to that of Udemy,
        providing a simple and clean UI/UX for them to interact with the service
        and classes. The Teacher experience is somewhat different in terms of
        signing up, but the actual course creation is very similar.
        \cite{skillshare-ux-1}\cite{skillshare-ux-2}
    }
    \item[LearnDash] {
        is a WordPress plugin that allows you to create and sell
        courses on your WordPress site. It is a very popular plugin, with
        over 50,000 active installs and a 4.5-star rating on the WordPress
        plugin repository. It is a very powerful plugin, with a lot of
        features, but it is also very complex and can be difficult to use.
        \cite{learndash}\\
        Being a plugin for WP, it is very easy to customise how the plugin
        looks by modifying the theme on the site to best suit the business /
        individual's needs. The plugin allows teachers to create courses and
        students to access them, with a lot of customisation options for the
        courses.
    }
    \item[Moodle] {
        is the closest comparison that I can find to the product that I was
        aiming to create. It is a PHP based LMS that is very powerful and
        feature-rich.
    }
\end{description}

\section{Tech Stacks}\label{sec:techstacks}

A `Tech Stack' is the term used to describe the combination of programming
languages, frameworks, libraries, and tools that are used to build a software
application. The choice of a tech stack is crucial as it can affect the
performance, scalability, and maintainability of the application. The following
are some of the popular tech stacks used in web development:

\begin{description}
    \item[LAMP/LEMP Stacks] {}
    \item[MEAN/MERN/MEVN Stacks] {}
\end{description}

\section{Databases and ORMs}\label{sec:databasesorms}

what is a database.
what is an ORM
popular databases and ORMs

%%%%%%%%%%%% Technical Specification %%%%%%%%%%%%

\chapter{Technical Specification}\label{ch:techspec}

\section{Framework}\label{sec:framework}

Which stack was chosen and why
which framework was chosen within that stack and why

\section{Database and ORM}\label{sec:databaseorm}

Which database was chosen and why
which ORM was chosen and why

%%%%%%%%%%%% Implementation %%%%%%%%%%%%

\chapter{Implementation}\label{ch:implementation}

\section{Authentication}\label{sec:authentication}

\section{User Management}\label{sec:usermanagement}

\section{CRUD Operations}\label{sec:crudoperations}

\section{Data Loading}\label{sec:dataloading}

\section{Known Issues}\label{sec:knownissues}

%%%%%%%%%%%% Testing %%%%%%%%%%%%

\chapter{Testing}\label{ch:testing}

\section{Unit Testing}\label{sec:unittesting}

\section{User Testing}\label{sec:usertesting}

%%%%%%%%%%%% Project Planning %%%%%%%%%%%%

\chapter {Project Planning}\label{ch:projectplanning}

%%%%%%%%%%%% Conclusion %%%%%%%%%%%%

\chapter{Conclusions}\label{ch:conclusion}

\bibliographystyle{ieeetr}
\bibliography{report}
\footnotetext{
    \textbf{Information}
    \emph{This document was written in LaTex and compiled with PDFTex and
        BibTex}
}

\end{document}
